% Auto-generated from diplom_asp/content/world_models.tex
\chapter{Использование моделей мира для оценки долгосрочной стоимости в безопасном обучении с подкреплением}
\label{ch:world_models}

Одним из основных ограничений безопасного обучения с подкреплением является необходимость принимать решения в условиях неполного знания о будущих последствиях действий агента. В классической постановке безопасного обучения ограничения задаются через функцию штрафа, а поведение агента оценивается не только по накопленному вознаграждению, но и по накопленной стоимости нарушений безопасности. Однако в процессе обучения агент не знает заранее, какие действия приведут к нарушению ограничений в будущем. Эта проблема особенно существенна в задачах с разреженными штрафами, где опасное действие может не приводить к немедленному нарушению, но переводить агента в область состояний, из которой дальнейшее безопасное поведение становится затруднительным.

Методы обучения с подкреплением на основе модели мира позволяют частично решить эту проблему за счёт предсказания последствий действий до их выполнения в реальной среде. В таких подходах обучаемая модель мира используется для оценки будущих состояний, вознаграждений и штрафов на коротком горизонте планирования. Однако оценка безопасности только на конечном горизонте планирования имеет принципиальное ограничение: она не учитывает стоимость нарушений, которые могут возникнуть после окончания планируемой траектории. Поэтому для безопасного исследования важно оценивать не только мгновенные штрафы, но и ожидаемую долгосрочную стоимость будущих нарушений.

В данной главе рассматривается метод SPOWL~\cite{spowl} как базовая архитектура безопасного обучения с подкреплением на основе модели мира. Основное внимание уделяется не всей архитектуре метода, а компонентам модели мира, отвечающим за оценку штрафов и долгосрочного риска. Исследуется способ построения функции долгосрочной стоимости штрафов, её использование при оценке кандидатных траекторий, а также способы агрегации ансамбля таких функций. Цель анализа состоит в том, чтобы определить, как выбор оценки стоимости влияет на баланс между достижением цели, исследованием среды и соблюдением ограничений безопасности.

\section{Оценка долгосрочной стоимости в модели мира SPOWL}
\label{sec:spowl_long_term_cost}

Метод SPOWL предназначен для безопасного обучения с подкреплением в непрерывных задачах управления. Он объединяет обучаемую модель мира, безопасную стратегию, планировщик безопасного улучшения и механизм выбора между действием стратегии и действием, предложенным планировщиком. В каждый момент времени агент может рассматривать две альтернативы: действие, сгенерированное безопасной стратегией, и действие, полученное с помощью планирования в модели мира. Итоговое действие выбирается на основе оценок полезности и стоимости.

В рамках данной главы общая архитектура SPOWL используется как вычислительный контекст для анализа оценки долгосрочной стоимости. Основным объектом рассмотрения является модель мира, которая предсказывает не только мгновенные вознаграждения и штрафы, но и функции стоимости для вознаграждения и штрафов. Наличие функции стоимости штрафов позволяет оценивать риск не только внутри горизонта планирования, но и за его пределами.

\subsection{Неявная модель мира}
\label{subsec:implicit_world_model}

Модель мира в SPOWL относится к классу неявных моделей мира. В отличие от моделей, которые явно восстанавливают наблюдения среды, неявная модель мира обучает латентное представление состояния и оптимизирует его относительно целевых величин, важных для принятия решений. Такой подход уменьшает необходимость в декодировании наблюдений и позволяет сосредоточить обучение модели на величинах, непосредственно используемых агентом: вознаграждении, штрафе, полезности и стоимости.

Пусть $s_t$ --- состояние среды, $a_t$ --- действие агента, а $z_t$ --- латентное представление состояния. Модель мира включает следующие компоненты:
\begin{equation}
    z_t = h(s_t),
\end{equation}
где $h$ --- кодировщик состояния;
\begin{equation}
    \hat{z}_{t+1} = f(z_t, a_t),
\end{equation}
где $f$ --- модель латентной динамики;
\begin{equation}
    \hat{r}_t = \hat{R}(z_t, a_t),
\end{equation}
где $\hat{R}$ --- модель мгновенного вознаграждения;
\begin{equation}
    \hat{Q}_{t,j} = \hat{Q}_j(z_t, a_t),
\end{equation}
где $\hat{Q}_j$ --- функции полезности вознаграждения;
\begin{equation}
    \hat{c}_{t,j} = \hat{C}_j(z_t, a_t),
\end{equation}
где $\hat{C}_j$ --- ансамбль моделей мгновенного штрафа;
\begin{equation}
    \hat{Q}^c_{t,j} = \hat{Q}^c_j(z_t, a_t),
\end{equation}
где $\hat{Q}^c_j$ --- ансамбль функций долгосрочной стоимости штрафов.

Ключевое отличие между $\hat{C}_j$ и $\hat{Q}^c_j$ состоит в том, что $\hat{C}_j$ оценивает локальный штраф для текущего перехода, тогда как $\hat{Q}^c_j$ оценивает ожидаемую накопленную стоимость будущих нарушений. Поэтому функция $\hat{Q}^c_j$ позволяет учитывать риск, который не проявляется в виде немедленного штрафа, но может возникнуть после нескольких последующих действий.

\subsection{Обучение функции долгосрочной стоимости штрафов}
\label{subsec:cost_value_training}

Функция долгосрочной стоимости штрафов обучается по принципу временных разностей. Для вознаграждения и штрафов используются отдельные целевые значения. Целевое значение для функции полезности вознаграждения имеет вид:
\begin{equation}
    Q_t = r_t + \gamma \overline{Q}(\hat{z}'_t, \hat{a}'_t),
\end{equation}
где $\gamma$ --- коэффициент дисконтирования вознаграждения, а $\overline{Q}$ --- целевая функция полезности.

Аналогично, целевое значение для функции стоимости штрафов задаётся как
\begin{equation}
    Q^c_t = c_t + \gamma_c \overline{Q}^c(\hat{z}'_t, \hat{a}'_t),
\end{equation}
где $c_t$ --- мгновенный штраф, $\gamma_c$ --- коэффициент дисконтирования стоимости, а $\overline{Q}^c$ --- целевая функция стоимости штрафов.

Эта формула показывает, что оценка стоимости действия состоит из двух частей: текущего штрафа и ожидаемой стоимости будущих нарушений. Следовательно, модель мира обучается не только обнаруживать непосредственные нарушения безопасности, но и предсказывать долгосрочный риск, связанный с текущим действием.

Полная функция потерь модели мира включает несколько компонент:
\begin{equation}
    \begin{aligned}
        L_t ={} &
        \|\hat{z}'_t - \operatorname{sg}(z'_t)\|^2_2
        + \operatorname{CE}(\hat{r}_t, r_t) \\
        &+ \frac{1}{N}\sum_{j=1}^{N}\operatorname{CE}(\hat{Q}_{t,j}, Q_t)
        + \frac{1}{N}\sum_{j=1}^{N}\operatorname{CE}(\hat{c}_{t,j}, c_t) \\
        &+ \frac{1}{N}\sum_{j=1}^{N}\operatorname{CE}(\hat{Q}^c_{t,j}, Q^c_t).
    \end{aligned}
\end{equation}
Здесь первая компонента отвечает за согласованность латентной динамики, вторая --- за предсказание вознаграждения, третья --- за обучение функций полезности, четвёртая --- за предсказание мгновенных штрафов, а последняя --- за обучение функций долгосрочной стоимости штрафов. Именно последняя компонента является центральной для оценки долгосрочного риска.

\subsection{Оценка стоимости кандидатной траектории}
\label{subsec:trajectory_cost_eval}

При планировании агент рассматривает кандидатную последовательность действий
\begin{equation}
    a_{t:t+H} = (a_t, a_{t+1}, \ldots, a_{t+H}),
\end{equation}
где $H$ --- горизонт планирования. Модель мира последовательно предсказывает латентные состояния и оценивает для каждой траектории ожидаемое вознаграждение и ожидаемую стоимость.

Схема использования этих оценок в процедуре планирования показана на рис.~\ref{fig:spowl_planning}. Планировщик формирует два множества кандидатных последовательностей действий: последовательности, порождённые текущей стратегией, и последовательности, полученные сэмплированием из нормального распределения. Для каждой последовательности модель мира вычисляет оценки $J^M$ и $J^M_c$. Затем на основе оценок траекторий, порождённых стратегией, задаются пороги $d^R$ и $d^c$, после чего выбираются элитные траектории, одновременно удовлетворяющие требованиям по вознаграждению и стоимости.

\begin{figure}[htbp]
    \centering
    \includegraphics[width=0.8\textwidth]{Images/plan.png}
    \caption{Схема безопасного планирования в SPOWL. Планировщик оценивает воображаемые траектории, полученные из текущей стратегии и из сэмплирования действий, вычисляет для них ожидаемое вознаграждение $J^M$ и стоимость $J^M_c$, после чего задаёт пороги $d^R$ и $d^c$ и выбирает элитные траектории.}
    \label{fig:spowl_planning}
\end{figure}

Для оценки вознаграждения используется сумма предсказанных мгновенных вознаграждений и остаточная оценка полезности в конце горизонта:
\begin{equation}
    J^M(z_t, a_{t:t+H}) =
    \sum_{i=0}^{H-1}\gamma^i \hat{r}_{t+i}
    + \gamma^H \hat{Q}_{\mathrm{agg}}(z_{t+H}, a_{t+H}).
\end{equation}

Для оценки безопасности аналогично используется сумма предсказанных мгновенных штрафов и функция долгосрочной стоимости штрафов:
\begin{equation}
    J^M_c(z_t, a_{t:t+H}) =
    \sum_{i=0}^{H-1}\gamma_c^i \hat{c}_{t+i}
    + \gamma_c^H \hat{Q}^c_{\mathrm{agg}}(z_{t+H}, a_{t+H}).
\end{equation}

Первая часть этой формулы отвечает за локальную стоимость траектории внутри горизонта планирования. Вторая часть отвечает за долгосрочный риск после окончания горизонта. Таким образом, функция $Q^c$ играет роль терминальной оценки стоимости, позволяющей учитывать потенциальные будущие нарушения, которые не видны при прямом моделировании только первых $H$ шагов.

Если использовать только сумму мгновенных штрафов,
\begin{equation}
    J^{M,\mathrm{local}}_c(z_t, a_{t:t+H}) =
    \sum_{i=0}^{H-1}\gamma_c^i \hat{c}_{t+i},
\end{equation}
то оценка безопасности становится локальной. Такой подход может быть достаточным, если штрафы плотные и нарушение быстро проявляется в пределах горизонта планирования. Однако в задачах с разреженными штрафами или отложенными последствиями локальная оценка может не увидеть будущий риск. Поэтому в SPOWL используется глобальная оценка стоимости, дополненная функцией долгосрочной стоимости штрафов:
\begin{equation}
    J^{M,\mathrm{global}}_c(z_t, a_{t:t+H}) =
    \sum_{i=0}^{H-1}\gamma_c^i \hat{c}_{t+i}
    + \gamma_c^H \hat{Q}^c_{\mathrm{agg}}(z_{t+H}, a_{t+H}).
\end{equation}

Таким образом, различие между локальной и глобальной оценкой состоит в наличии терминального слагаемого $\hat{Q}^c_{\mathrm{agg}}$. Это слагаемое позволяет учитывать риск за пределами короткого горизонта планирования.

\subsection{Агрегация ансамбля функций долгосрочной стоимости}
\label{subsec:ensemble_aggregation}

Поскольку функции стоимости аппроксимируются нейронными сетями, их оценки могут быть смещёнными. Недооценка стоимости приводит к выбору небезопасных действий, а переоценка стоимости приводит к чрезмерно консервативному поведению. Для повышения устойчивости оценки в модели мира используется ансамбль функций стоимости штрафов:
\begin{equation}
    \{\hat{Q}^c_1, \hat{Q}^c_2, \ldots, \hat{Q}^c_N\}.
\end{equation}

Итоговая оценка стоимости получается агрегацией значений ансамбля:
\begin{equation}
    \hat{Q}^c_{\mathrm{agg}}(z_t,a_t)
    =
    \operatorname{Agg}_{j=1}^{N}\hat{Q}^c_j(z_t,a_t).
\end{equation}

В данной работе рассматриваются три способа агрегации:
\begin{equation}
    \hat{Q}^c_{\min}(z_t,a_t)
    =
    \min_{j=1,\ldots,N}\hat{Q}^c_j(z_t,a_t),
\end{equation}
\begin{equation}
    \hat{Q}^c_{\max}(z_t,a_t)
    =
    \max_{j=1,\ldots,N}\hat{Q}^c_j(z_t,a_t),
\end{equation}
\begin{equation}
    \hat{Q}^c_{\mathrm{avg}}(z_t,a_t)
    =
    \frac{1}{N}\sum_{j=1}^{N}\hat{Q}^c_j(z_t,a_t).
\end{equation}

Минимальная агрегация соответствует оптимистичной оценке риска. Она выбирает наименьшую оценку стоимости среди элементов ансамбля и поэтому может недооценивать вероятность будущих нарушений. Такая оценка может способствовать исследованию среды, но повышает риск выбора небезопасных траекторий.

Максимальная агрегация соответствует пессимистичной оценке риска. Она выбирает наибольшую оценку стоимости и поэтому уменьшает вероятность выбора опасных действий. Однако такая оценка может приводить к чрезмерной консервативности, при которой агент избегает потенциально полезных траекторий и хуже достигает цели.

Усреднение по ансамблю является промежуточным вариантом. Оно снижает влияние отдельных выбросов в оценках функций стоимости и позволяет получить более устойчивую оценку риска. В задачах безопасного исследования такой баланс важен, поскольку агент должен одновременно избегать опасных действий и сохранять способность к исследованию среды.

\section{Экспериментальная проверка}
\label{sec:spowl_experiments}

В данном разделе рассматривается экспериментальная проверка роли оценки долгосрочной стоимости в безопасном обучении с подкреплением на основе модели мира. Основное внимание уделяется тому, как способ построения и агрегации функций стоимости штрафов влияет на безопасность и эффективность агента. При этом сначала приводится общее сравнение полной конфигурации SPOWL с базовыми методами, а затем рассматриваются абляционные эксперименты, связанные с оценкой долгосрочной стоимости.

Эксперименты направлены на ответы на следующие исследовательские вопросы.

\textbf{RQ1:} обеспечивает ли полная конфигурация SPOWL конкурентоспособное качество по сравнению с существующими методами безопасного обучения с подкреплением?

\textbf{RQ2:} как использование функции долгосрочной стоимости штрафов $Q^c$ влияет на оценку безопасности кандидатных траекторий по сравнению с локальной оценкой только по мгновенным штрафам?

\textbf{RQ3:} как способ агрегации ансамбля функций долгосрочной стоимости штрафов влияет на баланс между достижением цели и соблюдением ограничений безопасности?

\textbf{RQ4:} как размер ансамбля и размер подвыборки при агрегации функций стоимости влияют на итоговую безопасность и эффективность агента?

\subsection{Описание сред}
\label{subsec:spowl_envs}

Эксперименты проводились в средах SafetyGymnasium~\cite{ji2023safetygymnasium}, предназначенных для оценки алгоритмов безопасного обучения с подкреплением. Данный набор сред включает задачи непрерывного управления с различными типами агентов, динамикой и ограничениями безопасности.

В экспериментах используются следующие типы задач:
\begin{itemize}
    \item Goal: агент должен достичь целевой области, избегая опасных зон и объектов;
    \item Push: агент должен переместить объект в заданную область, не нарушая ограничения безопасности;
    \item Button: агент должен активировать целевую кнопку, взаимодействуя с объектами среды и избегая опасных зон.
\end{itemize}

Также рассматриваются разные типы тел агента:
\begin{itemize}
    \item Point --- простой точечный агент с низкоразмерной динамикой;
    \item Car --- агент с кинематическими ограничениями движения;
    \item Doggo --- многозвенный робот с более сложной динамикой;
    \item Ant --- многозвенный робот с высокой размерностью пространства состояний и действий.
\end{itemize}

В таблицах далее используются следующие сокращения: $P$ --- Point, $C$ --- Car, $D$ --- Doggo, $A$ --- Ant. Обозначения G1, G2, P1, B1 соответствуют задачам Goal1, Goal2, Push1 и Button1 соответственно. Например, PG1 обозначает задачу PointGoal1, а AG1 --- задачу AntGoal1.

Все методы обучались при одинаковом числе взаимодействий со средой. Итоговые значения метрик приводятся после $5 \cdot 10^6$ шагов взаимодействия.

\subsection{Метрики}
\label{subsec:spowl_metrics}

Для оценки методов используются три основные метрики.

\textbf{Return} --- суммарное вознаграждение за эпизод:
\begin{equation}
    R_{\mathrm{ep}} = \sum_{t=0}^{T_{\mathrm{ep}}} r_t.
\end{equation}
Эта метрика характеризует способность агента выполнять целевую задачу. Чем выше значение Return, тем эффективнее агент достигает цели.

\textbf{Costs} --- суммарная стоимость нарушений безопасности за эпизод:
\begin{equation}
    C_{\mathrm{ep}} = \sum_{t=0}^{T_{\mathrm{ep}}} c_t.
\end{equation}
Эта метрика показывает количество нарушений ограничений безопасности во время выполнения задачи. Чем ниже значение Costs, тем безопаснее поведение агента.

\textbf{Cost Rate} --- средний штраф за шаг во время обучения:
\begin{equation}
    \rho = \frac{1}{T}\sum_{t=0}^{T} c_t,
\end{equation}
где $T$ --- общее число шагов взаимодействия со средой. Эта метрика характеризует безопасность исследования в процессе обучения, а не только безопасность итоговой стратегии.

Результаты приводятся в виде среднего значения и стандартного отклонения по трём независимым запускам.

\subsection{Базовые методы}
\label{subsec:spowl_baselines}

Для сравнения используются как методы обучения с подкреплением на основе модели мира, так и безмодельные методы безопасного обучения с подкреплением.

\textbf{TD-MPC2}~\cite{hansen2024tdmpc2} используется как метод обучения с подкреплением на основе модели мира без учёта ограничений безопасности. Этот метод позволяет оценить, какой уровень целевой эффективности может быть достигнут при игнорировании штрафов. Поэтому TD-MPC2 обычно получает высокое вознаграждение, но накапливает большое число нарушений безопасности.

\textbf{SafeDreamer}~\cite{safedreamer} является методом безопасного обучения с подкреплением на основе модели мира на базе архитектуры Dreamer. Метод использует модель мира и планирование с учётом ограничений безопасности. Он является наиболее близким базовым методом обучения с подкреплением на основе модели мира, поскольку также использует предсказания модели мира для безопасного поведения.

\textbf{PPOLag}~\cite{stooke2020pidlag} --- безмодельный on-policy метод, основанный на PPO и лагранжевом учёте ограничений безопасности. Рассматриваются две версии метода: PPOLag(2) и PPOLag(25), где число в скобках обозначает допустимый порог стоимости за эпизод.

\textbf{CPO}~\cite{achiam2017cpo} --- on-policy метод оптимизации стратегии с ограничениями, в котором обновление стратегии строится с учётом ограничения на ожидаемую стоимость. В экспериментах используется версия CPO(25) с порогом стоимости 25.

\textbf{CVPO}~\cite{liu2022cvpo} --- off-policy метод безопасного обучения с подкреплением, использующий ограниченную оптимизацию стратегии. В экспериментах используется версия CVPO(25).

Для корректного сравнения все методы используют векторные наблюдения, включающие проприоцептивные признаки и информацию о ближайших объектах среды.

\subsection{Сравнение с базовыми методами (RQ1)}
\label{subsec:spowl_baseline_comparison}

Сначала проводится общее сравнение полной конфигурации SPOWL с базовыми методами. Это сравнение необходимо для проверки того, что предложенная конфигурация модели мира, включающая оценку долгосрочной стоимости штрафов и ансамблевую агрегацию функций стоимости, приводит не только к улучшению отдельных внутренних компонент, но и к конкурентоспособному итоговому поведению агента.

Результаты приведены в табл.~\ref{tab:spowl_baselines}. Во всех задачах оцениваются три метрики: $R$ --- эпизодическое вознаграждение, $C$ --- эпизодическая стоимость нарушений безопасности, $\rho$ --- частота стоимости во время обучения.

\begin{table}[ht]
    \caption{Сравнение SPOWL с базовыми методами на задачах SafetyGymnasium. TD-MPC2 не учитывает ограничения безопасности. Обозначения PPOLag(2), PPOLag(25), CPO(25), CVPO(25) указывают порог допустимой стоимости за эпизод. Метрика $R$ обозначает эпизодическое вознаграждение, $C$ --- эпизодическую стоимость, $\rho$ --- частоту стоимости во время обучения. Все значения приведены после $5 \cdot 10^6$ шагов взаимодействия со средой.}
    \label{tab:spowl_baselines}
    \centering
    \small
    \resizebox{\textwidth}{!}{
    \begin{tabular}{lr|llll|lll}
    \toprule
    Task & Metric & \textbf{SPOWL} & TD-MPC2 & SafeDreamer & PPOLag(2) & PPOLag(25) & CPO(25) & CVPO(25) \\
    \midrule
    \multirow{3}{*}{PG1}
        & $R\uparrow$ & $26.9 \pm 0.4$ & $27.5 \pm 0.2$ & $12.0 \pm 4.0$ & $8.5 \pm 3.5$ & $14.5 \pm 0.6$ & $19.2 \pm 3.0$ & $15.3 \pm 6.4$ \\
        & $C\downarrow$ & $1.5 \pm 1.1$ & $52.3 \pm 8.6$ & $10.3 \pm 9.5$ & $13.5 \pm 7.2$ & $19.7 \pm 3.0$ & $31.0 \pm 5.8$ & $22.9 \pm 8.5$ \\
        & $\rho \times 10^4\downarrow$ & $133 \pm 2$ & $515 \pm 9$ & $160 \pm 16$ & $222 \pm 19$ & $361 \pm 20$ & $275 \pm 21$ & $296 \pm 18$ \\
    \midrule
    \multirow{3}{*}{PG2}
        & $R\uparrow$ & $13.7 \pm 3.2$ & $27.4 \pm 0.5$ & $5.2 \pm 3.9$ & $1.8 \pm 1.7$ & $2.0 \pm 0.9$ & $3.2 \pm 3.2$ & $2.0 \pm 1.2$ \\
        & $C\downarrow$ & $4.8 \pm 1.2$ & $179.0 \pm 12.2$ & $0 \pm 0$ & $35.9 \pm 14.3$ & $33.7 \pm 16.3$ & $26.1 \pm 0.8$ & $22.6 \pm 21.3$ \\
        & $\rho \times 10^4\downarrow$ & $280 \pm 5$ & $1641 \pm 18$ & $221 \pm 32$ & $397 \pm 53$ & $570 \pm 73$ & $354 \pm 45$ & $391 \pm 6$ \\
    \midrule
    \multirow{3}{*}{PP1}
        & $R\uparrow$ & $18.7 \pm 1.4$ & $20.8 \pm 0.6$ & $12.5 \pm 7.2$ & $1.5 \pm 0.6$ & $4.4 \pm 2.7$ & $2.4 \pm 1.8$ & $10.4 \pm 2.7$ \\
        & $C\downarrow$ & $1.8 \pm 2.5$ & $34.2 \pm 11.6$ & $7.3 \pm 10.4$ & $33.5 \pm 16.4$ & $24.2 \pm 12.4$ & $25.5 \pm 8.1$ & $25.8 \pm 9.8$ \\
        & $\rho \times 10^4\downarrow$ & $103 \pm 5$ & $353 \pm 6$ & $71 \pm 3$ & $123 \pm 13$ & $276 \pm 16$ & $254 \pm 7$ & $253 \pm 12$ \\
    \midrule
    \multirow{3}{*}{PB1}
        & $R\uparrow$ & $10.0 \pm 1.5$ & $33.4 \pm 0.4$ & $0.1 \pm 1.4$ & $0.2 \pm 0.4$ & $3.7 \pm 2.5$ & $1.3 \pm 0.6$ & $10.1 \pm 1.7$ \\
        & $C\downarrow$ & $5.4 \pm 2.7$ & $103.3 \pm 3.1$ & $0.3 \pm 0.5$ & $43.7 \pm 35.4$ & $33.6 \pm 7.7$ & $22.8 \pm 5.4$ & $53.6 \pm 26.5$ \\
        & $\rho \times 10^4\downarrow$ & $269 \pm 3$ & $1219 \pm 10$ & $159 \pm 3$ & $395 \pm 142$ & $392 \pm 27$ & $336 \pm 12$ & $358 \pm 19$ \\
    \midrule
    \multirow{3}{*}{CG1}
        & $R\uparrow$ & $31.3 \pm 0.8$ & $34.8 \pm 0.3$ & $6.6 \pm 1.2$ & $3.5 \pm 3.2$ & $18.4 \pm 3.9$ & $21.6 \pm 2.4$ & $20.2 \pm 2.1$ \\
        & $C\downarrow$ & $1.0 \pm 1.0$ & $51.2 \pm 7.9$ & $0 \pm 0$ & $7.5 \pm 6.8$ & $30.5 \pm 5.0$ & $25.6 \pm 2.8$ & $34.8 \pm 3.4$ \\
        & $\rho \times 10^4\downarrow$ & $113 \pm 2$ & $563 \pm 11$ & $127 \pm 15$ & $214 \pm 6$ & $345 \pm 9$ & $321 \pm 10$ & $283 \pm 10$ \\
    \midrule
    \multirow{3}{*}{DG1}
        & $R\uparrow$ & $19.0 \pm 1.9$ & $27.9 \pm 1.5$ & $1.1 \pm 0.9$ & $0 \pm 0$ & $0 \pm 0$ & $0 \pm 0$ & $0 \pm 0$ \\
        & $C\downarrow$ & $0.3 \pm 0.4$ & $72.6 \pm 20.6$ & $0 \pm 0$ & $0 \pm 0$ & $21.8 \pm 30.8$ & $0 \pm 0$ & $0 \pm 0$ \\
        & $\rho \times 10^4\downarrow$ & $286 \pm 15$ & $659 \pm 33$ & $260 \pm 70$ & $249 \pm 15$ & $376 \pm 70$ & $224 \pm 79$ & $332 \pm 24$ \\
    \midrule
    \multirow{3}{*}{AG1}
        & $R\uparrow$ & $38.9 \pm 1.4$ & $61.0 \pm 2.6$ & $0.2 \pm 0.2$ & $0.2 \pm 0.2$ & $5.1 \pm 3.9$ & $0 \pm 0$ & $4.3 \pm 1.9$ \\
        & $C\downarrow$ & $2.0 \pm 0.2$ & $70.9 \pm 6.0$ & $0 \pm 0$ & $0 \pm 0$ & $23.6 \pm 18.2$ & $0 \pm 0$ & $41.5 \pm 44.1$ \\
        & $\rho \times 10^4\downarrow$ & $275 \pm 6$ & $682 \pm 6$ & $218 \pm 58$ & $220 \pm 141$ & $638 \pm 115$ & $466 \pm 38$ & $598 \pm 33$ \\
    \bottomrule
    \end{tabular}
    }
\end{table}

Из табл.~\ref{tab:spowl_baselines} видно, что TD-MPC2 достигает высокой целевой эффективности, но систематически нарушает ограничения безопасности, поскольку не учитывает штрафы при обучении и планировании. Например, в задачах PG2, PB1 и AG1 он получает высокое вознаграждение, но накапливает значительно большую стоимость нарушений по сравнению со SPOWL.

SafeDreamer, напротив, часто обеспечивает низкую стоимость, но во многих задачах делает это за счёт сильного ухудшения целевой эффективности. В задачах DoggoGoal1 и AntGoal1 значения вознаграждения близки к нулю, что указывает на неспособность метода решить задачу при сохранении безопасности.

Безмодельные методы PPOLag и CPO также не обеспечивают устойчивого баланса между вознаграждением и стоимостью. Уменьшение допустимого порога стоимости в PPOLag(2) снижает исследовательскую активность и ухудшает целевую эффективность, тогда как более мягкий порог PPOLag(25) улучшает достижение цели, но часто приводит к росту нарушений безопасности.

CVPO показывает более высокую эффективность по сравнению с on-policy методами в некоторых задачах, однако во многих случаях достигает этого ценой большего числа нарушений безопасности. Например, в задаче PB1 CVPO получает сравнимое со SPOWL вознаграждение, но имеет существенно большую эпизодическую стоимость.

Таким образом, общее сравнение показывает, что полная конфигурация SPOWL обеспечивает более устойчивый компромисс между достижением цели и соблюдением ограничений безопасности. В отличие от TD-MPC2, метод учитывает стоимость нарушений; в отличие от чрезмерно консервативных безопасных методов, он сохраняет ненулевую целевую эффективность в сложных задачах. Этот результат подтверждает, что компоненты модели мира, отвечающие за оценку долгосрочной стоимости и её ансамблевую агрегацию, являются практически значимыми для итогового качества безопасного поведения агента.

\subsection{Сравнение локальной и глобальной оценки стоимости (RQ2)}
\label{subsec:spowl_local_global}

Первый эксперимент направлен на проверку роли функции долгосрочной стоимости штрафов $Q^c$. Для этого сравниваются два способа оценки стоимости кандидатной траектории при планировании: локальная оценка, учитывающая только предсказанные мгновенные штрафы на конечном горизонте, и глобальная оценка, дополнительно использующая функцию долгосрочной стоимости штрафов.

Локальная оценка использует только сумму предсказанных мгновенных штрафов на горизонте планирования:
\begin{equation}
\label{eq:spowl_local_cost}
    J^{M,\mathrm{local}}_c =
    \sum_{i=0}^{H-1}\gamma_c^i \hat{c}_{t+i}.
\end{equation}

Глобальная оценка дополнительно использует терминальную оценку долгосрочной стоимости штрафов:
\begin{equation}
\label{eq:spowl_global_cost}
    J^{M,\mathrm{global}}_c =
    \sum_{i=0}^{H-1}\gamma_c^i \hat{c}_{t+i}
    + \gamma_c^H \hat{Q}^c_{\mathrm{avg}}(z_{t+H}, a_{t+H}).
\end{equation}

Таким образом, локальная оценка проверяет безопасность только внутри короткого горизонта планирования, тогда как глобальная оценка учитывает возможные нарушения после завершения этого горизонта за счёт функции $Q^c$.

Для сравнения использовался планировщик с фиксированным порогом допустимой стоимости $d_{\mathrm{plan}}$. В случае глобальной оценки стоимость траектории вычислялась по формуле~\eqref{eq:spowl_global_cost}, то есть с учётом терминальной оценки $\hat{Q}^c_{\mathrm{avg}}$. Результаты приведены на рис.~\ref{fig:spowl_global_estimation}.

\begin{figure}[htbp]
    \centering
    \includegraphics[width=0.95\textwidth]{Images/cceg.png}
    \caption{Сравнение фиксированных порогов безопасности при глобальной оценке стоимости траектории. Глобальная оценка использует сумму предсказанных мгновенных штрафов и терминальную оценку долгосрочной стоимости $\hat{Q}^c_{\mathrm{avg}}$.}
    \label{fig:spowl_global_estimation}
\end{figure}

В случае локальной оценки стоимость траектории вычислялась по формуле~\eqref{eq:spowl_local_cost}, то есть только как сумма предсказанных мгновенных штрафов на конечном горизонте планирования. Результаты приведены на рис.~\ref{fig:spowl_local_estimation}.

\begin{figure}[htbp]
    \centering
    \includegraphics[width=0.95\textwidth]{Images/ccel.png}
    \caption{Сравнение фиксированных порогов безопасности при локальной оценке стоимости траектории. Локальная оценка использует только сумму предсказанных мгновенных штрафов на горизонте планирования и не учитывает терминальную оценку $\hat{Q}^c_{\mathrm{avg}}$.}
    \label{fig:spowl_local_estimation}
\end{figure}

Результаты на рис.~\ref{fig:spowl_global_estimation} и~\ref{fig:spowl_local_estimation} показывают, что локальная и глобальная оценки стоимости приводят к различному поведению агента при выборе порога безопасности. Локальная оценка обладает ограниченной видимостью будущего. В средах с разреженными штрафами нарушение может не проявиться в пределах горизонта планирования, поэтому агент получает слишком оптимистичную оценку безопасности. Это может ускорять рост вознаграждения на ранних этапах обучения, но одновременно повышает риск небезопасного поведения. Чтобы поддерживать безопасность при локальной оценке, необходимо задавать очень малые пороги допустимой стоимости, что ухудшает достижение цели.

Глобальная оценка ведёт себя иначе. Поскольку она использует функцию долгосрочной стоимости $Q^c$, на ранних этапах обучения оценка риска может быть высокой. Это ограничивает исследование среды, но позволяет учитывать будущие нарушения, которые не видны при локальной оценке. По мере обучения модели и стратегии оценка стоимости становится более точной, после чего агент получает возможность выбирать более продуктивные траектории без резкого роста нарушений безопасности.

Дополнительно на рис.~\ref{fig:spowl_thresholds} показана динамика глобальной оценки стоимости при разных фиксированных порогах. На ранних этапах обучения средняя оценка долгосрочной стоимости остаётся высокой, поэтому планирование ограничивает выбор рискованных траекторий. Рост эффективности начинается после того, как оценка стоимости становится ниже заданного порога.

\begin{figure}[htbp]
    \centering
    \includegraphics[width=\linewidth]{Images/thresholds.png}
    \caption{Динамика глобальной оценки стоимости при фиксированных порогах безопасности. На ранних этапах обучения средняя оценка долгосрочной стоимости высока, поэтому планирование ограничивает исследование. Рост эффективности начинается после того, как оценка стоимости становится ниже заданного порога.}
    \label{fig:spowl_thresholds}
\end{figure}

Следовательно, данный эксперимент показывает, что функция $Q^c$ является необходимым компонентом для оценки долгосрочного риска в безопасном обучении с подкреплением на основе модели мира. Локальная сумма мгновенных штрафов недостаточна для устойчивого безопасного исследования, поскольку не учитывает стоимость за пределами короткого горизонта планирования. Глобальная оценка, включающая $\hat{Q}^c_{\mathrm{avg}}$, позволяет учитывать отложенный риск и тем самым даёт более полную оценку безопасности кандидатной траектории.

\subsection{Сравнение способов агрегации ансамбля функций стоимости (RQ3)}
\label{subsec:aggregation_comparison}

Второй эксперимент исследует, как способ агрегации ансамбля функций стоимости штрафов влияет на безопасность и эффективность агента. Рассматриваются три варианта агрегации: минимум, максимум и среднее.

Минимальная агрегация задаёт оптимистичную оценку риска:
\begin{equation}
    \hat{Q}^c_{\min}(z_t,a_t)
    =
    \min_j \hat{Q}^c_j(z_t,a_t).
\end{equation}

Этот вариант может быть полезен для уменьшения консервативности, но в безопасном обучении он опасен, поскольку даже одна заниженная оценка в ансамбле может привести к выбору небезопасной траектории. Экспериментально такая агрегация приводит к недооценке стоимости и росту числа нарушений безопасности.

Максимальная агрегация задаёт пессимистичную оценку риска:
\begin{equation}
    \hat{Q}^c_{\max}(z_t,a_t)
    =
    \max_j \hat{Q}^c_j(z_t,a_t).
\end{equation}

Она уменьшает вероятность выбора опасных действий, но может переоценивать стоимость и тем самым снижать качество решения задачи. В результате агент становится более консервативным, хуже исследует среду и может получать меньшее суммарное вознаграждение.

Усреднение задаёт сбалансированную оценку:
\begin{equation}
    \hat{Q}^c_{\mathrm{avg}}(z_t,a_t)
    =
    \frac{1}{N}\sum_{j=1}^{N}\hat{Q}^c_j(z_t,a_t).
\end{equation}

Экспериментальные результаты показывают, что усреднение по ансамблю позволяет избежать крайних эффектов. Оно не приводит к такой сильной недооценке риска, как минимум, и не создаёт чрезмерной консервативности, характерной для максимума. Поэтому средняя агрегация обеспечивает наиболее устойчивый баланс между безопасностью и достижением цели. Сравнение способов агрегации приведено на рис.~\ref{fig:cvaltarg}.

\begin{figure}[htbp]
    \centering
    \includegraphics[width=0.85\linewidth]{Images/cvaltarg.png}
    \caption{Сравнение способов агрегации ансамбля функций долгосрочной стоимости штрафов.}
    \label{fig:cvaltarg}
\end{figure}

Этот результат показывает, что сама по себе функция долгосрочной стоимости недостаточна: важно также корректно агрегировать оценки ансамбля. Неправильная агрегация может изменить поведение агента даже при той же архитектуре модели мира и тех же данных обучения.

\subsection{Влияние размера ансамбля и размера подвыборки (RQ4)}
\label{subsec:ensemble_size_ablation}

Третий эксперимент направлен на анализ параметров ансамбля функций стоимости. Исследуется влияние двух факторов: общего размера ансамбля и размера подвыборки, используемой при агрегации оценок.

Если ансамбль слишком мал, оценка стоимости становится нестабильной и сильно зависит от ошибок отдельной функции. В этом случае агент может либо недооценивать риск, либо переоценивать его в зависимости от случайного смещения одной модели. Эксперименты показывают, что использование единственной функции стоимости уступает ансамблевым вариантам по совокупности метрик.

Увеличение размера ансамбля повышает устойчивость оценки, однако чрезмерное увеличение не даёт существенного улучшения. Это связано с тем, что после определённого размера ансамбля дополнительное усреднение уже слабо изменяет итоговую оценку, но увеличивает вычислительные затраты.

Размер подвыборки также влияет на безопасность. Более крупные подвыборки позволяют получать более устойчивую оценку стоимости и уменьшают число нарушений безопасности. Однако слишком пессимистичная агрегация, например максимум по крупной подвыборке, может снижать вознаграждение без соответствующего выигрыша по безопасности. В таких случаях агент становится более осторожным, но не получает существенного уменьшения нарушений. Результаты абляционного исследования приведены на рис.~\ref{fig:ensembles}.

\begin{figure}[htbp]
    \centering
    \includegraphics[width=\linewidth]{Images/ensembles.png}
    \caption{Абляционное исследование размера ансамбля и подвыборки функций долгосрочной стоимости штрафов.}
    \label{fig:ensembles}
\end{figure}

На основании проведённого анализа в SPOWL используется схема $5 \times 5$ с усреднением: пять функций в ансамбле и агрегация средним значением. Такая конфигурация обеспечивает достаточно устойчивую оценку стоимости без заметного ухудшения целевой эффективности.

\subsection{Итоговый анализ экспериментальных результатов}
\label{subsec:experiments_summary}

Результаты экспериментов позволяют выделить несколько важных свойств оценки долгосрочной стоимости.

Во-первых, безопасность в обучении на основе модели мира не может надёжно оцениваться только по мгновенным штрафам внутри короткого горизонта планирования. Такая оценка является локальной и может пропускать будущие нарушения, особенно в средах с разреженными штрафами.

Во-вторых, функция долгосрочной стоимости штрафов $Q^c$ позволяет учитывать отложенный риск. Она дополняет сумму мгновенных штрафов терминальной оценкой будущей стоимости и тем самым расширяет горизонт оценки безопасности без необходимости явно моделировать длинные траектории.

В-третьих, ансамблевая оценка необходима для повышения устойчивости. Отдельная функция стоимости может давать смещённые оценки, что в безопасном обучении приводит либо к нарушениям ограничений, либо к чрезмерной консервативности.

В-четвёртых, способ агрегации ансамбля существенно влияет на поведение агента. Минимальная агрегация приводит к недооценке риска, максимальная --- к переоценке риска, а усреднение обеспечивает более сбалансированную оценку.

Таким образом, основной результат данной главы состоит в том, что устойчивое безопасное исследование в обучении с подкреплением на основе модели мира требует не только предсказания мгновенных штрафов, но и корректной оценки долгосрочной стоимости будущих нарушений. Эта оценка должна строиться с использованием ансамбля функций стоимости и сбалансированного способа агрегации.

\section{Выводы по главе}
\label{sec:chapter2_conclusion}

В данной главе рассмотрена задача оценки долгосрочной стоимости нарушений безопасности в обучении с подкреплением на основе модели мира. В качестве базовой архитектуры использован метод SPOWL, однако основное внимание было уделено не всей архитектуре метода, а компонентам модели мира, отвечающим за предсказание мгновенных штрафов и оценку ожидаемой накопленной стоимости будущих нарушений.

Показано, что при безопасном планировании в модели мира недостаточно учитывать только сумму мгновенных штрафов на конечном горизонте. Такая локальная оценка стоимости ограничена длиной горизонта планирования и может не учитывать нарушения, возникающие после завершения кандидатной траектории. Поэтому в модели мира используется функция долгосрочной стоимости штрафов $Q^c$, обучаемая по принципу временных разностей. Эта функция позволяет дополнить локальную сумму предсказанных штрафов терминальной оценкой будущего риска.

В рамках метода SPOWL функция $Q^c$ используется при вычислении стоимости кандидатных траекторий в планировщике. Стоимость траектории строится как сумма предсказанных мгновенных штрафов и остаточной долгосрочной стоимости в конце горизонта. Такое представление позволяет учитывать не только непосредственные нарушения безопасности, но и потенциально опасные продолжения траектории за пределами короткого горизонта планирования.

Экспериментальное сравнение локальной и глобальной оценки стоимости показывает, что локальная оценка может ускорять исследование за счёт более оптимистичной оценки риска, однако в задачах с разреженными штрафами она обладает ограниченной видимостью будущих нарушений. Для поддержания безопасности при такой оценке требуются малые пороги допустимой стоимости, что ухудшает достижение цели. Глобальная оценка, включающая $Q^c$, позволяет учитывать долгосрочный риск, хотя на ранних этапах обучения может приводить к более осторожному поведению из-за высоких оценок стоимости.

Также рассмотрена агрегация ансамбля функций долгосрочной стоимости штрафов. Эксперименты показывают, что способ агрегации существенно влияет на баланс между безопасностью и эффективностью. Минимальная агрегация соответствует оптимистичной оценке риска и может приводить к недооценке стоимости, что повышает число нарушений безопасности. Максимальная агрегация, напротив, переоценивает риск и может снижать целевую эффективность агента из-за чрезмерной консервативности. Усреднение по ансамблю обеспечивает более устойчивый баланс между достижением цели и соблюдением ограничений безопасности.

Дополнительный анализ размера ансамбля и размера подвыборки показывает, что ансамблевая оценка повышает устойчивость функции стоимости по сравнению с одиночной оценкой. При этом чрезмерное увеличение ансамбля не даёт существенного улучшения, а слишком консервативная агрегация может ухудшать вознаграждение без сопоставимого выигрыша по безопасности. Наиболее сбалансированным вариантом является усреднение по ансамблю, обеспечивающее достаточно точную оценку стоимости без заметной деградации целевой эффективности.

Общее сравнение с базовыми методами показывает, что полная конфигурация SPOWL обеспечивает устойчивый компромисс между достижением цели и соблюдением ограничений безопасности. В отличие от TD-MPC2, метод учитывает стоимость нарушений и существенно снижает число небезопасных действий. В отличие от чрезмерно консервативных безопасных методов, SPOWL сохраняет ненулевую целевую эффективность в более сложных задачах. Это подтверждает практическую значимость корректной оценки долгосрочной стоимости в модели мира.

Таким образом, результаты главы показывают, что устойчивое безопасное исследование в обучении с подкреплением на основе модели мира требует оценки не только мгновенных штрафов, но и ожидаемой накопленной стоимости будущих нарушений. Корректное построение и агрегация ансамбля функций долгосрочной стоимости являются ключевыми условиями для баланса между безопасностью и эффективностью агента: недооценка риска приводит к нарушениям ограничений, а переоценка риска --- к консервативному поведению и снижению вознаграждения.
